\subsection{Decomposition into extremal elements}

Let A be an integral domain, K its field of fractions and U the multiplicative group of invertible elements of A. Recall (Algebra, Chapter VI, § 1, no. 5) that there is a canonical isomorphism of \(\mathbf{K}^* / \mathbf{U}\) onto the group \(\mathcal{P}^*\) of non-zero fractional principal ideals of A. Condition (b) of Theorem 1 may then be translated as follows:

PROPOSITION 2. Let A be an integral domain. For A to be factorial, it is necessary and sufficient that there exist a subset P of A such that every \(a \in A - \{0\}\) may be written uniquely in the form \(a = u \prod_{p \in P} p^{n(p)}\) , where \(u \in U\) and the \(n(p)\) are positive integers which are zero except for a finite number of them.

If P satisfies this condition, clearly all its elements are extremal and every extremal element of A is associated with a unique element of P. Recall that P is then called a representative system of extremal elements of A (Algebra, Chapter VII, § 1, no. 3, Definition 2).

Suppose always that A is factorial. It has been seen (no. 2, Theorem 1) that the group \(\mathcal{P}^*\) is a lattice. We may therefore apply the results of Algebra, Chapter VI, §1, nos. 9 and 13. In particular, every element of \(\mathbf{K}^*\) may be written in an essentially unique way in the form of an irreducible fraction. Any
two elements \(a, b\) of \(\mathbf{K}^*\) have a g.c.d. and a l.c.m.; if \(a = u\prod_{p\in \mathbb{P}}p^{n(p)}\) and

\[
b = u ^ {\prime} \prod_ {p \in P} p ^ {m (p)}
\]

are decompositions of \(a\) and \(b\) as products of extremal elements, then:

\begin{enumerate}
\def\labelenumi{(\arabic{enumi})}
\tightlist
\item
\end{enumerate}

\[
\mathbf {g . c . d .} (a, b) = w \prod_ {p \in P} p ^ {\inf (m (p), n (p))}\tag{2}
\]

\[
\mathbf {l . c . m .} (a, b) = w ^ {\prime} \prod_ {p \in P} p ^ {\sup (m (p), n (p))}
\]

where \(w, w'\) are in U. We recover, in particular, the results of Algebra, Chapter VII, § 1, no. 3.

For all \(p \in P\) , the mapping \(a \mapsto n(p)\) is a discrete valuation \(v_{p}\) on K whose ring is obviously \(A_{Ap}\) . It follows from Theorem 1(e) that the \(v_{p}\) are just the essential valuations of A and that the ideals \(A_{p}\) ( \(p \in P\) ) are just the prime ideals of A of height 1.

